\section{Materials and Experimental Setup}
\label{sec:materials}

\subsection{Dataset}
\label{sec:dataset}

This study was conducted using a publicly available brain tumor Magnetic
Resonance Imaging (MRI) dataset containing three tumor categories: Glioma,
Meningioma, and Pituitary Tumor~\cite{cheng2017}. The dataset comprises
$3{,}064$ T1-weighted contrast-enhanced MRI slices collected from clinical
examinations and has been widely used as a benchmark in brain tumor
classification research.

% NOTE: notebook uses the FULL dataset (MAX_SAMPLES_PER_CLASS = None), i.e.
% 3064 images, NOT a balanced 1500 subset. Confirm per-class counts before
% submission if the notebook is re-run with subsampling.
The class distribution is naturally imbalanced, comprising approximately
$1{,}426$ glioma, $708$ meningioma, and $930$ pituitary images (imbalance ratio
$\approx 2.0$). Rather than artificially resampling the data, the native
distribution is preserved and class imbalance is addressed through stratified
data partitioning, class-weighted training where supported by the classifier,
and macro-averaged evaluation metrics that assign equal importance to every
class. Figure~\ref{fig:samples} illustrates representative examples from the
three classes.

\begin{figure}[t]
\centering
\includegraphics[width=\linewidth]{dataset}
\caption{Representative MRI slices from the three tumor classes
(glioma, meningioma, and pituitary tumor).}
\label{fig:samples}
\end{figure}

The classification task is formulated as a three-class supervised learning
problem, with class labels $0$, $1$, and $2$ corresponding to glioma,
meningioma, and pituitary tumor, respectively.

\subsection{Image Preprocessing}
\label{sec:preprocessing}

Medical images often exhibit variations in intensity distribution, image
quality, and acquisition conditions. Therefore, a preprocessing pipeline was
applied prior to feature extraction to enhance image consistency and improve
descriptor robustness.

\subsubsection{Grayscale Conversion}
Since the MRI images contain limited color information and most handcrafted
descriptors operate on intensity values, all images were represented in
grayscale. This reduces computational complexity while preserving
diagnostically relevant structural information.

\subsubsection{Image Resizing}
To ensure uniform feature extraction across the dataset, all images were resized
to a fixed spatial resolution of $128 \times 128$ pixels using anti-aliased
interpolation. Standardizing image dimensions prevents descriptor
inconsistencies caused by varying image sizes and facilitates fair comparison
between samples.

\subsubsection{Contrast Enhancement}
Contrast Limited Adaptive Histogram Equalization (CLAHE) was applied
(clip limit $0.02$) to improve local contrast and enhance subtle tumor-related
structures. Unlike conventional histogram equalization, CLAHE limits noise
amplification while increasing the visibility of clinically relevant texture
patterns.

\subsubsection{Intensity Normalization}
Pixel intensity values were normalized to the range $[0,1]$ to reduce variations
introduced by different acquisition settings and imaging devices, improving the
stability and reproducibility of the extracted feature representations.

The preprocessing pipeline can therefore be summarized as: MRI image
$\rightarrow$ grayscale $\rightarrow$ resizing $\rightarrow$ CLAHE enhancement
$\rightarrow$ intensity normalization (Figure~\ref{fig:preprocessing}).

\begin{figure}[t]
\centering
\includegraphics[width=\linewidth]{preprocessing_pipeline}
\caption{Preprocessing pipeline applied to each MRI slice before feature
extraction.}
\label{fig:preprocessing}
\end{figure}

\subsection{Overall Experimental Framework}
\label{sec:framework}

The proposed framework consists of two sequential phases designed to separate
feature optimization from classifier evaluation.

\paragraph{Phase 1: Separability-Guided Feature Configuration}
The first phase identifies the most discriminative parameter configurations for
each handcrafted descriptor family. Candidate configurations are evaluated using
classifier-independent separability metrics, allowing feature quality to be
assessed without introducing bias from any particular learning algorithm. The
descriptor families investigated are the Gray-Level Co-occurrence Matrix (GLCM),
Local Binary Patterns (LBP), the Discrete Wavelet Transform (DWT), and the
Histogram of Oriented Gradients (HOG). For each descriptor, a search space of
parameter combinations is explored and ranked according to a composite
separability score.

\paragraph{Phase 2: Classifier Benchmarking}
The optimized descriptors obtained in Phase~1 are used to construct multiple
feature representations, which are then evaluated using a diverse set of machine
learning classifiers under identical experimental conditions. This two-phase
design ensures that feature optimization and classifier benchmarking remain
independent, enabling a more objective assessment of both components.

\subsection{Feature Representation and Cleaning}
\label{sec:feat_design}

From the optimized descriptors, eight feature representations were constructed:
the individual families (A: first-order statistics; B: GLCM+LBP texture;
C: DWT; D: HOG) and four fusions (A+B, B+C, A+B+C, and the complete set
Full(opt)). Each representation is passed through a cleaning and standardization
pipeline (variance filtering, correlation pruning, and training-set z-score
normalization) before classification. The construction of these representations
and the cleaning pipeline are detailed in Section~\ref{sec:features}.

\subsection{Machine Learning Classifiers}
\label{sec:classifiers_setup}

Eight machine learning algorithms spanning linear, distance-based, tree-based,
ensemble, boosting, and neural-network paradigms were evaluated: SVM, KNN,
Random Forest, Extra Trees, XGBoost, LightGBM, Logistic Regression, and a
Multi-Layer Perceptron. Each algorithm and its hyperparameter search space is
described in Section~\ref{sec:phase2}.

\subsection{Hyperparameter Optimization}
\label{sec:hyperparams_setup}

To ensure a fair comparison between classifiers, hyperparameter optimization was
performed using grid search combined with stratified cross-validation. For each
classifier, a predefined search space was explored, and the configuration
achieving the highest macro-averaged $F_1$ score in cross-validation was
selected. Macro-$F_1$ was adopted as the selection criterion so that model
tuning is not dominated by the majority class. This procedure reduces the risk
of underestimating classifier performance due to suboptimal parameter choices.

\subsection{Evaluation Protocol}
\label{sec:protocol}

The dataset was divided into training and testing subsets using a stratified
hold-out strategy ($80\%/20\%$) that preserves class proportions. Within the
training set, model selection and hyperparameter tuning were performed using
five-fold stratified cross-validation. Performance was reported using accuracy,
together with macro-averaged precision, recall, and $F_1$ score. Macro-averaging
ensures equal importance across all tumor classes, which is appropriate given
the imbalanced class distribution.

\subsection{Statistical Validation}
\label{sec:stats_setup}

Beyond conventional performance metrics, statistical validation was conducted to
assess the significance and robustness of observed performance differences. The
following analyses were employed: Cohen's kappa coefficient~\cite{cohen1960};
bootstrap $95\%$ confidence intervals; McNemar's test~\cite{mcnemar1947} for
pairwise comparison of classifiers on the same test set; the Wilcoxon
signed-rank test; and the Friedman rank test with Nemenyi post-hoc
analysis~\cite{demsar2006} for comparison across multiple
classifiers and feature sets. These procedures provide a rigorous framework for
determining whether performance improvements are statistically meaningful rather
than the result of random variation.

The methodology described in this section establishes a reproducible
experimental framework for evaluating separability-guided handcrafted feature
optimization and machine learning classifier benchmarking in brain tumor MRI
classification.
